\SourceLeaf{122}
\Article{28}
In dieser Abhandlung sind wir der allgemein üblichen Erklärung der magnetischen Erscheinungen gefolgt, weil sie diesen Erscheinungen vollständig gerecht wird und mit weit einfacheren Rechnungen auskommt als die Auffassung, die den Magnetismus galvano-elektrischen Kreisen um die Teilchen des magnetischen Körpers zuschreibt. Unser Urteil über diese Auffassung, die sich gewiss in mehrfacher Hinsicht empfiehlt, wollten wir offenlassen. Sie zu bestätigen oder zu verwerfen wäre verfrüht gewesen, da das Gesetz der gegenseitigen Wirkung zwischen den Elementen solcher Kreise noch nicht hinreichend erforscht zu sein scheint. Welche Auffassung von den rein magnetischen und den elektromagnetischen Erscheinungen man künftig auch annimmt, für die rein magnetischen Erscheinungen muss sie überall zum gleichen Ergebnis führen wie die übliche Auffassung. Was wir in dieser Abhandlung auf dieser Grundlage entwickelt haben, kann sich allein in der Form ändern, sein Gehalt bleibt derselbe.
